% GENERATED FILE. Edit canonical lessons, then run npm run generate:books.
% canonical-source-hash: fnv1a64:f0fe8e36d5ba313a
% canonical-lessons: PT-C525-bibliotecario, PT-C525-requisitar, PT-C525-editora, PT-C525-traducao, PT-C525-tradutor

\chapter{Na biblioteca}
\label{ch:pt-na-biblioteca}
\hlchaptermodality{525}

\begin{chapteropening}
\textbf{By the end of this chapter:} \emph{I can talk about borrowing books, publishers and translation.}

Complete the last lesson of chapter 525: I can talk about borrowing books, publishers and translation.
\end{chapteropening}

\section[o bibliotecário]{o bibliotecário — librarian}

\label{lesson:PT-C525-bibliotecario}



\begin{quote}
Before the new one: say the Portuguese for rhyme, then the Portuguese for to recite.
\end{quote}

\subsection*{You'll want to know: o bibliotecário}
\textbf{o bibliotecário} — \textquotedblleft{}librarian\textquotedblright{}.

\begin{quote}
O bibliotecário ajudou-me a encontrar o livro.
\end{quote}

\begin{quote}
\emph{The librarian helped me find the book.}
\end{quote}

A woman is \emph{a bibliotecária}.

\begin{cousinweb}
\emph{bibliotecário} is built on \emph{biblioteca}, a learned borrowing through Latin \emph{bibliothēca} \textquotedblleft{}a library\textquotedblright{} (Lewis \& Short) from Greek \emph{biblioth\'{\={e}}kē}, from \emph{biblíon} \textquotedblleft{}book\textquotedblright{} and \emph{th\'{\={e}}kē} \textquotedblleft{}a store\textquotedblright{}.
\end{cousinweb}

\begin{grammarlens}[title={What you've built}]
One more word for telling, explaining and describing at B1.
\end{grammarlens}

\subsection*{Your turn}
\begin{itemize}
\raggedright
  \item \emph{Say it:} \emph{o bibliotecário}
  \item \emph{Say it:} \emph{o bibliotecário}, once more
  \item \emph{Say it:} a sentence of your own with \emph{o bibliotecário}
  \item \emph{Recall:} say the Portuguese for rhyme, then the Portuguese for to recite, then say \emph{o bibliotecário} again
\end{itemize}

\begin{tcolorbox}[breakable,colback=teal!4,colframe=teal!35!black,title={Before you move on}]
What does o bibliotecário mean? (\textquotedblleft{}Librarian\textquotedblright{}.) Say it once more.
\end{tcolorbox}
\section[requisitar]{requisitar — to borrow (from a library); to request}

\label{lesson:PT-C525-requisitar}



\begin{quote}
Before the new one: say the Portuguese for to recite, then the Portuguese for librarian.
\end{quote}

\subsection*{You'll want to know: requisitar}
\textbf{requisitar} — \textquotedblleft{}to borrow (from a library); to request\textquotedblright{}.

\begin{quote}
Requisitei três livros na biblioteca.
\end{quote}

\begin{quote}
\emph{I borrowed three books from the library.}
\end{quote}

In Portugal the usual verb for library loans.

\begin{cousinweb}
\emph{requisitar} is built on Latin \emph{requīsītus} \textquotedblleft{}sought, asked for\textquotedblright{} (Lewis \& Short), from \emph{requīrere}, built on \emph{quaerere} \textquotedblleft{}to seek\textquotedblright{}.
\end{cousinweb}

\begin{grammarlens}[title={What you've built}]
One more word for telling, explaining and describing at B1.
\end{grammarlens}

\subsection*{Your turn}
\begin{itemize}
\raggedright
  \item \emph{Say it:} \emph{requisitar}
  \item \emph{Say it:} \emph{requisitar}, once more
  \item \emph{Say it:} a sentence of your own with \emph{requisitar}
  \item \emph{Recall:} say the Portuguese for to recite, then the Portuguese for librarian, then say \emph{requisitar} again
\end{itemize}

\begin{tcolorbox}[breakable,colback=teal!4,colframe=teal!35!black,title={Before you move on}]
What does requisitar mean? (\textquotedblleft{}To borrow (from a library); to request\textquotedblright{}.) Say it once more.
\end{tcolorbox}
\section[a editora]{a editora — publisher, publishing house}

\label{lesson:PT-C525-editora}



\begin{quote}
Before the new one: say the Portuguese for librarian, then the Portuguese for to borrow.
\end{quote}

\subsection*{You'll want to know: a editora}
\textbf{a editora} — \textquotedblleft{}publisher, publishing house\textquotedblright{}.

\begin{quote}
A editora vai publicar o livro em maio.
\end{quote}

\begin{quote}
\emph{The publisher is going to bring out the book in May.}
\end{quote}

Built on \emph{editar} \textquotedblleft{}to publish\textquotedblright{}.

\begin{cousinweb}
\emph{editora} is built on \emph{editar} \textquotedblleft{}to publish\textquotedblright{}, which came through French \emph{éditer} (EtymDB) from Latin \emph{ēdere} \textquotedblleft{}to give out, to put forth\textquotedblright{} (Lewis \& Short).
\end{cousinweb}

\begin{grammarlens}[title={What you've built}]
One more word for telling, explaining and describing at B1.
\end{grammarlens}

\subsection*{Your turn}
\begin{itemize}
\raggedright
  \item \emph{Say it:} \emph{a editora}
  \item \emph{Say it:} \emph{a editora}, once more
  \item \emph{Say it:} a sentence of your own with \emph{a editora}
  \item \emph{Recall:} say the Portuguese for librarian, then the Portuguese for to borrow, then say \emph{a editora} again
\end{itemize}

\begin{tcolorbox}[breakable,colback=teal!4,colframe=teal!35!black,title={Before you move on}]
What does a editora mean? (\textquotedblleft{}Publisher, publishing house\textquotedblright{}.) Say it once more.
\end{tcolorbox}
\section[a tradução]{a tradução — translation}

\label{lesson:PT-C525-traducao}



\begin{quote}
Before the new one: say the Portuguese for to borrow, then the Portuguese for publisher.
\end{quote}

\subsection*{You'll want to know: a tradução}
\textbf{a tradução} — \textquotedblleft{}translation\textquotedblright{}.

\begin{quote}
Li o romance numa boa tradução.
\end{quote}

\begin{quote}
\emph{I read the novel in a good translation.}
\end{quote}

\emph{Tradução simultânea} — simultaneous interpreting.

\begin{cousinweb}
\emph{tradução} comes from Latin \emph{trāductiō} \textquotedblleft{}a leading across\textquotedblright{} (EtymDB), from \emph{trādūcere}, built on \emph{trāns} \textquotedblleft{}across\textquotedblright{} and \emph{dūcere} \textquotedblleft{}to lead\textquotedblright{}.
\end{cousinweb}

\begin{grammarlens}[title={What you've built}]
One more word for telling, explaining and describing at B1.
\end{grammarlens}

\subsection*{Your turn}
\begin{itemize}
\raggedright
  \item \emph{Say it:} \emph{a tradução}
  \item \emph{Say it:} \emph{a tradução}, once more
  \item \emph{Say it:} a sentence of your own with \emph{a tradução}
  \item \emph{Recall:} say the Portuguese for to borrow, then the Portuguese for publisher, then say \emph{a tradução} again
\end{itemize}

\begin{tcolorbox}[breakable,colback=teal!4,colframe=teal!35!black,title={Before you move on}]
What does a tradução mean? (\textquotedblleft{}Translation\textquotedblright{}.) Say it once more.
\end{tcolorbox}
\section[o tradutor]{o tradutor — translator}

\label{lesson:PT-C525-tradutor}



\begin{quote}
Before the new one: say the Portuguese for publisher, then the Portuguese for translation.
\end{quote}

\subsection*{You'll want to know: o tradutor}
\textbf{o tradutor} — \textquotedblleft{}translator\textquotedblright{}.

\begin{quote}
O tradutor trabalha do inglês para o português.
\end{quote}

\begin{quote}
\emph{The translator works from English into Portuguese.}
\end{quote}

A woman is \emph{a tradutora}.

\begin{cousinweb}
\emph{tradutor} is built like \emph{tradução} on Latin \emph{trādūcere} \textquotedblleft{}to lead across\textquotedblright{} (Lewis \& Short).
\end{cousinweb}

\begin{grammarlens}[title={What you've built}]
One more word for telling, explaining and describing at B1.
\end{grammarlens}

\subsection*{Your turn}
\begin{itemize}
\raggedright
  \item \emph{Say it:} \emph{o tradutor}
  \item \emph{Say it:} \emph{o tradutor}, once more
  \item \emph{Say it:} a sentence of your own with \emph{o tradutor}
  \item \emph{Recall:} say the Portuguese for publisher, then the Portuguese for translation, then say \emph{o tradutor} again
\end{itemize}

\begin{tcolorbox}[breakable,colback=teal!4,colframe=teal!35!black,title={Before you move on}]
What does o tradutor mean? (\textquotedblleft{}Translator\textquotedblright{}.) Say it once more.
\end{tcolorbox}
